%%%%%%%%%%%%%%%%%%%%%%%%%%%%%%%%%%%%%%%%%%%%%%%%%%%%%%%%%%%%%%%%%%
%% Three problems from the Erdos corpus
%%%%%%%%%%%%%%%%%%%%%%%%%%%%%%%%%%%%%%%%%%%%%%%%%%%%%%%%%%%%%%%%%%
\documentclass[11pt]{article}

\usepackage[margin=1in]{geometry}
\usepackage{amsmath,amssymb,amsthm}
\usepackage{mathtools}
\usepackage{hyperref}
\usepackage{booktabs}
\usepackage{microtype}

\theoremstyle{plain}
\newtheorem{theorem}{Theorem}[section]
\newtheorem{lemma}[theorem]{Lemma}
\newtheorem{proposition}[theorem]{Proposition}
\newtheorem{corollary}[theorem]{Corollary}
\newtheorem{conjecture}[theorem]{Conjecture}

\theoremstyle{definition}
\newtheorem{definition}[theorem]{Definition}
\newtheorem{problem}[theorem]{Open Problem}
\newtheorem{remark}[theorem]{Remark}

\newcommand{\DS}{\mathrm{DS}}
\newcommand{\Rcop}{R_{\mathrm{cop}}}
\newcommand{\Gcop}{G_{\mathrm{cop}}}

\begin{document}

\title{Three new problems from the Erd\H{o}s corpus}

\author{Alexander Towell\thanks{Department of Computer Science,
Southern Illinois University Edwardsville.
M.S.\ Mathematics, M.S.\ Computer Science; Ph.D.\ student (Computer Science).
ORCID: 0000-0001-6443-9897.
Email: \texttt{atowell@siue.edu}}}

\date{}
\maketitle

\begin{abstract}
I introduce three problems arising from a systematic study of the
Erd\H{o}s problem corpus. (1)~\emph{Density-relaxed Schur numbers}
$\DS(k,\alpha)$ interpolate between Schur numbers and density Ramsey
theory; we determine $\DS(2,\alpha)$ for all~$\alpha$, revealing a
sharp phase transition at $\alpha=1/2$, with a complete Lean~4
formalization.
(2)~\emph{Coprime Ramsey numbers} $\Rcop(k)$ measure Ramsey forcing
on the coprime graph; we prove $\Rcop(3)=11$ by exhaustive extension.
(3)~I show that the set of primes, which maximizes $\sum 1/(a\log a)$
over primitive sets (Lichtman~\cite{lichtman2023}), is \emph{not}
optimal for coprime pair counts, and give a lower bound
$M^*(n)\ge\frac{64}{9\pi^2}\binom{\lceil n/2\rceil}{2}(1+o(1))$.
\end{abstract}

\smallskip\noindent\textbf{MSC 2020:} 05D10, 05C55, 11B75, 11N25.

%% =================================================================
\section{Introduction}\label{sec:intro}
%% =================================================================

A systematic study of the 1{,}135 problems in the Erd\H{o}s problem
database~\cite{tao2025erdos} reveals natural questions at the
intersection of Erd\H{o}s's interests---Schur numbers, coprime graphs,
primitive sets---that do not appear to have been formulated. I present
three such problems with solutions and computational results. All code
and Lean formalizations are publicly available.\footnote{Code and data:
\texttt{doi.org/10.5281/zenodo.18638635}}

%% =================================================================
\section{Density-relaxed Schur numbers}\label{sec:ds}
%% =================================================================

A \emph{Schur triple} is a triple $(a,b,c)$ with $a+b=c$ and
$a,b\ge 1$ (so $a=b$ is permitted).

\begin{definition}\label{def:ds}
For $k\ge 1$ and $\alpha\in(0,1]$, a $k$-coloring of $[N]$ is
\emph{$(k,\alpha)$-avoiding} if every color class is both sum-free
and has at most $\alpha N$ elements. Define
\[
\DS(k,\alpha) \;=\; 1+\max\{N\ge 1: \text{a $(k,\alpha)$-avoiding
coloring of $[N]$ exists}\},
\]
with $\DS(k,\alpha)=1$ if no such $N$ exists.
\end{definition}

When $\alpha=1$ the density constraint is vacuous, so
$\DS(k,1)=S(k)+1$ where $S(k)$ is the $k$-th Schur
number~\cite{schur1916}. When $\alpha<1/k$, pigeonhole forces some
class to exceed $\alpha N$ elements for all $N\ge 1$, so
$\DS(k,\alpha)=1$. Since Schur triples force at $S(k)+1$
independently of density, we always have $\DS(k,\alpha)\le S(k)+1$.

\begin{theorem}\label{thm:ds}
$\DS(2,\alpha)=5$ for all $\alpha\ge 1/2$, and $\DS(2,\alpha)=1$ for
$\alpha<1/2$. In particular $\DS(2,\tfrac{1}{2})=5=S(2)+1$.
\end{theorem}

\begin{proof}
For $\alpha<1/2$: in any $2$-coloring of $[N]$, some class has
$\lceil N/2\rceil>\alpha N$ elements (since $\alpha<1/2$), so no
$(2,\alpha)$-avoiding coloring exists for any~$N$.

For $\alpha\ge 1/2$: the coloring $C_1=\{1,4\}$, $C_2=\{2,3\}$ of
$[4]$ has both classes sum-free with $|C_i|=2\le\alpha\cdot 4$
(using $\alpha\ge 1/2$). So $N=4$ is $(2,\alpha)$-avoiding.
Since $S(2)=4$, every $2$-coloring of $[5]$ contains a Schur
triple, so $N=5$ is not avoiding. Hence $\DS(2,\alpha)=1+4=5$.
\end{proof}

\begin{remark}\label{rem:transition}
The density-forcing threshold ($\alpha=1/2=1/k$ for $k=2$)
coincides with the Schur-forcing threshold ($S(2)+1=5$),
collapsing $\DS(2,\cdot)$ to a step function. For $k\ge 3$, these
thresholds may separate, yielding a non-trivial interpolation regime.
\end{remark}

\begin{remark}\label{rem:lean}
Theorem~\ref{thm:ds} has been fully formalized in Lean~4 (Mathlib),
with 12 theorems proved and 0 \texttt{sorry} statements. The
formalization is in \texttt{lean/NPG2\_DensitySchur.lean}.
\end{remark}

\begin{problem}\label{prob:ds}
Determine $\DS(k,\alpha)$ for $k\ge 3$ and $\alpha\in(1/k,1)$.
Is there a non-trivial interpolation regime, i.e., does
$\DS(k,\alpha)<S(k)+1$ for some $\alpha>1/k$?
\end{problem}

%% =================================================================
\section{Coprime Ramsey numbers}\label{sec:rcop}
%% =================================================================

\begin{definition}\label{def:rcop}
The \emph{coprime graph} $\Gcop(n)$ has vertex set $[n]$ with
$\{i,j\}$ an edge iff $\gcd(i,j)=1$. Define $\Rcop(k)$ as the
minimum $n$ such that every $2$-coloring of $E(\Gcop(n))$ contains a
monochromatic $K_k$.
\end{definition}

Since $E(\Gcop(n))\subseteq E(K_n)$, any $2$-coloring of $K_n$
restricts to a $2$-coloring of $\Gcop(n)$; if the former avoids
monochromatic $K_k$, so does the latter (any monochromatic $K_k$ in
$\Gcop(n)$ would use only coprime edges, hence appear in $K_n$ too).
Therefore $\Rcop(k)\ge R(k,k)$~\cite{ramsey1930}.

The clique number satisfies $\omega(\Gcop(n))=\pi(n)+1$: the set
$\{1\}\cup\{\text{primes}\le n\}$ is a clique of this size. For the
upper bound, note that in any pairwise coprime set $S\subseteq\{2,
\ldots,n\}$, each element's prime factorization uses a disjoint set
of primes (else two elements share a factor); since each element
uses $\ge 1$ prime and all primes used are $\le n$, we have
$|S|\le\pi(n)$.

\begin{theorem}\label{thm:rcop3}
$\Rcop(3)=11$.
\end{theorem}

\begin{proof}
By exhaustive extension from $n=8$.

\emph{Base.} Among all $2^{21}$ edge $2$-colorings of $\Gcop(8)$
($21$ edges), exactly $36$ avoid monochromatic triangles.

\emph{Extension.} For $n=9$: vertex~$9$ adds $6$ coprime edges;
extending $36$ colorings through $2^6$ options each yields $36$
avoiding. For $n=10$: vertex~$10$ adds $4$ edges; extension yields
$156$ avoiding. For $n=11$: vertex~$11$ is prime, adding $10$ edges
(coprime to all of $1,\ldots,10$); extending $156$ colorings through
$2^{10}$ options each yields $0$ avoiding.

\begin{center}
\begin{tabular}{@{}cccc@{}}
\toprule
$n$ & Edges in $\Gcop(n)$ & New edges & Avoiding colorings \\
\midrule
8  & 21 & ---  & 36  \\
9  & 27 & 6  & 36  \\
10 & 31 & 4  & 156 \\
11 & 41 & 10 & 0   \\
\bottomrule
\end{tabular}
\end{center}
\end{proof}

The increase at $n=10$ reflects that $10=2\times 5$ has only $4$
coprime neighbors, imposing weak constraints. The collapse at $n=11$
is driven by $11$ being prime: its $10$ edges create too many
triangles for any $2$-coloring to avoid.

\begin{remark}\label{rem:rcop-ratio}
$\Rcop(3)/R(3,3)=11/6\approx 1.83$ exceeds the inverse coprime
density $\pi^2/6\approx 1.64$. If $\Gcop(n)$ behaved like a random
graph of density $6/\pi^2$, one would predict $\Rcop(3)\approx 10$.
The excess arises from multiplicative correlations among missing
edges (multiples of a prime~$p$ form correlated
non-adjacencies)~\cite{erdos1997coprime}.
\end{remark}

\begin{conjecture}\label{conj:rcop4}
$\Rcop(4)=20$. \emph{Evidence:} An avoiding coloring at $n=19$ was
found by random search ($\sim\!160{,}000$ trials); no avoiding
coloring was found at any $n\in\{20,\ldots,49\}$ ($\sim\!10^6$
random trials each). This is heuristic evidence only---the extension
method is infeasible at this scale ($\sim\!1.47\times 10^6$ base
colorings).
\end{conjecture}

\begin{problem}\label{prob:rcop}
Determine the growth rate of $\Rcop(k)$. Is $\Rcop(k)>R(k,k)$ for
all $k\ge 3$?
\end{problem}

%% =================================================================
\section{Coprime pairs in primitive sets}\label{sec:prim}
%% =================================================================

A set $A\subseteq[n]$ is \emph{primitive} if no element divides
another~\cite{erdos1935}. By Dilworth's theorem~\cite{dilworth1950},
every primitive $A\subseteq[n]$ has $|A|\le\lceil n/2\rceil$.
Let $M(A)=|\{(a,b)\in A^2:a<b,\;\gcd(a,b)=1\}|$ and
$M^*(n)=\max\{M(A):A\subseteq[n]\text{ primitive}\}$.

Lichtman~\cite{lichtman2023} proved that primes maximize
$f(A)=\sum_{a\in A}1/(a\log a)$ over primitive $A$---the Erd\H{o}s
primitive set conjecture. We show primes do \emph{not} maximize $M$.

The \emph{top layer} $T(n)=\{\lfloor n/2\rfloor+1,\ldots,n\}$ is
primitive with $|T(n)|=\lceil n/2\rceil$ (Dilworth-maximum).
Define the \emph{shifted top layer} $S(n)$ by replacing each even
$2k\in T(n)$ with $k$ when $k$ is odd and $k>n/3$.

\begin{lemma}\label{lem:prim}
$S(n)$ is primitive with $|S(n)|=\lceil n/2\rceil$.
\end{lemma}

\begin{proof}
Replacement is bijective (each $2k$ maps to a distinct $k\in(n/3,
n/2]\setminus T(n)$), preserving size. For primitivity: retained
elements of $T(n)$ exceed $n/2$, so no two divide each other. A new
element $k\in(n/3,n/2]$ cannot divide a retained element $b>n/2$:
the multiple $2k$ was removed, and $3k>n\ge b$. Two new elements
$k_1<k_2$ with $k_1\mid k_2$ would require $k_2\ge 2k_1>2n/3>
n/2\ge k_2$, a contradiction.
\end{proof}

\begin{proposition}\label{prop:density}
As $n\to\infty$, $M(T(n))\sim\frac{6}{\pi^2}\binom{\lceil
n/2\rceil}{2}$ and $M(S(n))\sim\frac{64}{9\pi^2}\binom{\lceil
n/2\rceil}{2}$.
\end{proposition}

\begin{proof}
For $T(n)$: the proportion of coprime pairs among $L$ consecutive
integers converges to $\prod_p(1-1/p^2)=6/\pi^2$ as $L\to\infty$
(classical; see, e.g., Hardy--Wright~\cite{hardywright}, Thm.~332). Since $T(n)$ consists
of $\lceil n/2\rceil$ consecutive integers, the claim follows.

For $S(n)$: the even fraction satisfies $f_E\to 1/3$ (we replace
roughly $1/6$ of the $\lceil n/2\rceil$ elements, reducing even
elements from $1/2$ to $1/3$). Two even elements share factor~$2$,
contributing $0$ coprime pairs. For pairs with at least one odd
element, coprimality is governed by odd primes; since the elements
of $S(n)$ are drawn from intervals of length $\Theta(n)$ in $[n]$,
the coprime probability for such pairs converges to
$\prod_{p\ge 3}(1-1/p^2)=8/\pi^2$. Weighting by parity classes:
$d(S(n))\to\frac{8}{\pi^2}(1-f_E^2)=\frac{64}{9\pi^2}$.
\end{proof}

\begin{theorem}\label{thm:main-prim}
$M^*(n)\ge\frac{64}{9\pi^2}\binom{\lceil n/2\rceil}{2}(1+o(1))$.
For $n\ge 10$, $M(P(n))<M^*(n)$ where $P(n)$ is the set of primes
up to $n$; the latter claim is verified by exhaustive computation.
\end{theorem}

\begin{proof}
The lower bound is realized by $S(n)$
(Lemma~\ref{lem:prim} and Proposition~\ref{prop:density}).
Since $M(P(n))=\binom{\pi(n)}{2}\sim n^2/(2\ln^2 n)=o(n^2)$ while
$M(S(n))=\Theta(n^2)$, primes are eventually suboptimal.
Exhaustive computation confirms this for $n\ge 10$:

\begin{center}
\begin{tabular}{@{}clcc@{}}
\toprule
$n$ & Optimal primitive set & $M^*(n)$ & $M(P(n))$ \\
\midrule
8   & $\{2,3,5,7\}$                    & 6  & 6 \\
10  & $\{4,5,6,7,9\}$                  & 8  & 6 \\
15  & $\{4,6,7,9,10,11,13,15\}$        & 21 & 15 \\
18  & $\{4,6,7,9,10,11,13,15,17\}$     & 29 & 21 \\
\bottomrule
\end{tabular}
\end{center}
\end{proof}

\begin{conjecture}\label{conj:prim-asymptotic}
$M^*(n)\sim\frac{64}{9\pi^2}\binom{\lceil n/2\rceil}{2}$.
That is, $S(n)$ is asymptotically optimal among primitive sets
for maximizing coprime pairs.
\end{conjecture}

\begin{remark}\label{rem:structure}
Optimal sets use prime powers ($4=2^2$, $9=3^2$) and products
($6=2\cdot 3$, $10=2\cdot 5$) instead of small primes.
Primes have coprime density $d=1$ but size only $\pi(n)$;
Dilworth-maximum sets have size $\lceil n/2\rceil$ but lower density.
Since $M\propto d\cdot|A|^2$, size dominates.
\end{remark}

\begin{problem}\label{prob:func}
Characterize which functionals $F$ on primitive sets are maximized by
primes. Lichtman's $\sum 1/(a\log a)$ favors primes; $M(A)$ does not.
\end{problem}

%% =================================================================
\section*{Acknowledgments}
%% =================================================================

The AI system Claude (Anthropic, model Claude Opus 4.6) was used
extensively in this work: it produced the computational experiments
(Python, 1{,}826 tests), constructed the proofs, implemented the
Lean~4 formalization, and drafted this manuscript. The author
directed the research program, selected problems, and reviewed all
mathematical claims. The proofs were additionally subjected to three
rounds of critical review and a Lean~4 formalization (0
\texttt{sorry} statements for Section~\ref{sec:ds}).

The Erd\H{o}s problem database~\cite{tao2025erdos} was essential to
this work. I thank Terry Tao and the contributors to
\texttt{github.com/teorth/erdosproblems} for creating and maintaining
this resource.

%% =================================================================
\begin{thebibliography}{9}

\bibitem{dilworth1950}
R.~P.~Dilworth.
A decomposition theorem for partially ordered sets.
\emph{Ann.\ of Math.\ (2)}, 51:161--166, 1950.

\bibitem{erdos1935}
P.~Erd\H{o}s.
Note on sequences of integers no one of which is divisible by any other.
\emph{J.\ London Math.\ Soc.}, 10:126--128, 1935.

\bibitem{hardywright}
G.~H.~Hardy and E.~M.~Wright.
\emph{An Introduction to the Theory of Numbers}.
Oxford University Press, 6th edition, 2008.

\bibitem{erdos1997coprime}
P.~Erd\H{o}s and A.~S\'ark\"ozy.
On cycles in the coprime graph of integers.
\emph{Electron.\ J.\ Combin.}, 4(2):R8, 1997.

\bibitem{lichtman2023}
J.~D.~Lichtman.
A proof of the {E}rd\H{o}s primitive set conjecture.
\emph{Forum Math.\ Pi}, 11:e18, 2023.

\bibitem{ramsey1930}
F.~P.~Ramsey.
On a problem of formal logic.
\emph{Proc.\ London Math.\ Soc.\ (2)}, 30:264--286, 1930.

\bibitem{schur1916}
I.~Schur.
\"{U}ber die {K}ongruenz $x^m+y^m\equiv z^m\pmod{p}$.
\emph{Jahresber.\ Deutsch.\ Math.-Verein.}, 25:114--117, 1916.

\bibitem{tao2025erdos}
T.~Tao et~al.
Erd\H{o}s problems database.
\texttt{github.com/teorth/erdosproblems}, 2025.

\end{thebibliography}

\end{document}
